% ============================================================
\chapter{共通問題（問5）の対策と試験本番の戦略}\label[appendix]{app:G}
% ============================================================
統計応用（医薬生物学）は大問5問から3問を選ぶ90分の記述式試験で，問5は人文科学・社会科学・理工学・医薬生物学の4分野に共通の問題である
（本書で確認した2016〜2025年の9回すべてでこの構成．問題冊子の注記による）．
本付録では，(1) 9回分の問5を年度別に整理し，(2) 必要な数理統計の知識と本書の対応を確認し，(3) 不足する知識を最小限補い，
(4) 問題の選び方・途中点・時間配分を実践的にまとめる．

\begin{warnbox}[この付録の読み方]
\begin{itemize}
\item 「問5は常に易しい」「問5は必ず選ぶべき」とは言えない．次の\cref{sec:G-cases}で見るように，問5を選んで成功した事例と失敗した事例の両方がある．
\item 難易度は年度と受験者の習熟度に依存するので，本付録では一律の難易度評価を付けない．代わりに「式変形の段数」「数値計算の量」「誘導の有無」「前の小問への依存」という\textbf{観察できる特徴}で記述する．
\item 過去9回の出題実績は，次の試験で何が出るかの確率ではない．問5は9回で同じテーマが繰り返されていない．
\end{itemize}
\end{warnbox}

\section{年度別の整理}\label{sec:G-years}
各年の問5を本書の分析で整理した（問題文は転載しない）．「本書の対応」は本書の章・付録で直接扱う内容があるかを示す．
\begin{center}\scriptsize
\setlength{\tabcolsep}{3pt}
\begin{longtable}{@{}cL{30mm}L{30mm}L{48mm}L{22mm}@{}}
\toprule
年 & テーマ & 分野 & 必要な数理統計の知識・導出技術 & 本書の対応\\
\midrule
\endfirsthead
\toprule
年 & テーマ & 分野 & 必要な数理統計の知識・導出技術 & 本書の対応\\
\midrule
\endhead
\bottomrule
\endfoot
2016 & 2標本 $t$ 検定と信頼区間の重なり & 推測（区間・検定） & $t$ 分布の CI，プールした分散，2標本 $t$，$t$ 分位点の大小と不等式の整理 & \cref{ch:01,ch:03}\\
2017 & 割付法と群サイズの分布，CI 幅から $n$ & 確率モデル，症例数 & 二項分布，経路の確率の列挙，分布の対称性，既知分散の CI 幅，二項の正規近似と $\sqrt n$ の2次方程式 & \cref{ch:02,ch:12}（逐次割付の確率計算はなし）\\
2018 & 正規2成分の混合と二峰性 & 確率分布 & 混合分布の平均・分散（全分散の公式），正規密度の2回微分，極値の条件 & \cref{subsec:G-mixture}\\
2019 & 適合度 $\chi^2$ 検定，ABO 血液型の EM & 分割表，最尤法 & 適合度 $\chi^2$ と自由度の根拠（多項→多変量正規），ラグランジュ乗数法，E ステップ & \cref{ch:16}，\cref{subsec:G-gof}\\
2021 & 診断検査の正規モデル & 確率（Bayes の定理） & 正規分布の標準化，全確率の公式，PPV・NPV と有病率 & \cref{ch:05}\\
2022 & 平均への回帰 & 2変量正規・回帰 & 差の分散，2変量正規の条件付き分布（与えられる），真値＋誤差モデル & \cref{sec:08-ancova,sec:17-normal}，\cref{subsec:G-rtm}\\
2023 & モンティ・ホール問題と二項検定 & 確率，検定 & Bayes の定理，正確な二項検定の棄却域と検出力 & \cref{ch:05}，\cref{subsec:G-exact}\\
2024 & 重回帰の抑制効果 & 回帰・線形モデル & 標準化変数の正規方程式（2×2 の逆行列），決定係数 & \cref{subsec:G-reg}\\
2025 & ランダム化回答法 & 標本調査，最尤法 & 全確率の公式，再パラメータ化した MLE，パラメータ空間の境界での MLE & \cref{subsec:G-rr}\\
\end{longtable}
\end{center}
2025年の解答は公式略解に導出を補ったもので，略解 PDF との照合は本書では行っていない．

\paragraph{横断的な観察}
\begin{itemize}
\item 分野は「確率（条件付き確率・Bayes の定理）」「確率モデル・分布」「推測（区間・検定・検出力・症例数）」が中心で，1年の中で2〜3分野を組み合わせる形が多い．
      9回の範囲では，時系列，ノンパラメトリック検定，主成分分析などの多変量解析手法は出ていない．
\item 題材（教授法，試験得点，血圧，採用試験，年収調査，ゲームなど）は分野固有の専門知識を要しない．必要な公式（条件付き分布など）は問題文で与えられることが多い．
\item デルタ法，積率母関数，尤度比検定，十分統計量は9回の問5には出ていない（医薬の問1〜4や統計数理では頻出）．
\end{itemize}

\section{必要な数理統計のスキルと本書の対応}\label{sec:G-skills}
9回のうちそのスキルを使う小問があった年の数（本書の分析）は，多い順に
正規・$t$・$\chi^2$ の標準化と数表の読み取り（6），解釈の記述（6），期待値・分散・共分散の計算（5），条件付き確率・Bayes の定理（5），
区間・検定・検出力・症例数の構成（4），二項・多項分布（4），最尤推定（2），2×2 の行列計算（2），微分による関数形状の解析（1）である．
9回中5回（2016，2017，2019，2021，2022）は初版の章で主要部分に対応できた．
残る4回（2018 混合分布，2023 正確な二項検定の検出力，2024 重回帰，2025 ランダム化回答法）と，2019年の適合度検定の根拠，2022年の平均への回帰の解釈を，次節で最小限補う．

\section{本文にない知識の最小限の補強}\label{sec:G-gaps}
いずれも統計数理の標準的な内容であり，ここでは問5で使う形に絞って要点と導出の筋を示す．

\subsection{適合度 \texorpdfstring{$\chi^2$}{χ²} 検定と自由度}\label{subsec:G-gof}
$K$ カテゴリの度数 $(O_1,\dots,O_K)\sim$ 多項分布 $(n;p_1,\dots,p_K)$ で，$H_0:p_i=p_{i0}$ を
\[
X^2=\sum_{i=1}^K\frac{(O_i-np_{i0})^2}{np_{i0}}\ \approx\ \chi^2_{K-1}\quad(H_0\text{ の下，}n\text{ が大きいとき})
\]
で検定する．\textbf{根拠}：$Z_i=(O_i-np_{i0})/\sqrt{np_{i0}}$ とすると，多項分布の共分散 $\Cov(O_i,O_j)=n p_i(\delta_{ij}-p_j)$ より
$\bm Z$ の共分散行列は $\bm I-\sqrt{\bm p}\sqrt{\bm p}\transp$（$\sqrt{\bm p}=(\sqrt{p_{10}},\dots,\sqrt{p_{K0}})\transp$）である．これは階数 $K-1$ の射影行列なので，
中心極限定理で $\bm Z$ が多変量正規に近づくと $X^2=\|\bm Z\|^2$ は $K-1$ 個の独立な標準正規の2乗和に近づく．
\textbf{自由度}：$\sum O_i=n$ という1つの制約で1減り，$p_{i0}$ を $q$ 個のパラメータの推定値で置き換えるとさらに $q$ 減って $K-1-q$ になる．
度数をすべて $k$ 倍すると $X^2$ も $k$ 倍になる（同じ割合でも標本が大きいほど有意になりやすい）．

\subsection{重回帰（2説明変数）と抑制効果}\label{subsec:G-reg}
すべての変数を標準化すると，$y$ を $x_1,x_2$ に回帰した係数 $\bm b$ は正規方程式 $\bm R_{xx}\bm b=\bm r_{xy}$ の解で
\[
b_1=\frac{r_{y1}-r_{12}r_{y2}}{1-r_{12}^2},\qquad b_2=\frac{r_{y2}-r_{12}r_{y1}}{1-r_{12}^2},\qquad R^2=\V[\hat y]=b_1r_{y1}+b_2r_{y2}
\]
（$r_{y1}=\Corr(y,x_1)$ など．$\bm R_{xx}=\begin{pmatrix}1&r_{12}\\r_{12}&1\end{pmatrix}$ の逆行列は $\frac1{1-r_{12}^2}\begin{pmatrix}1&-r_{12}\\-r_{12}&1\end{pmatrix}$）．
$R^2$ は $\V[\hat y]=\bm b\transp\bm R_{xx}\bm b=\bm b\transp\bm r_{xy}$ と，$\Cov(y,\hat y)=\bm b\transp\bm r_{xy}$ から得られる．
\textbf{抑制効果}：$y$ とほぼ無相関な $x_2$ でも，$x_1$ と相関していれば $b_2\ne0$ となり，$x_1$ の係数を $r_{y1}$ より大きくして $R^2$ を $r_{y1}^2$ より大きくする（$x_2$ が $x_1$ の中の $y$ と無関係な成分を取り除く）．
例：$r_{y1}=0.5$，$r_{y2}=0.1$，$r_{12}=0.6$ なら $b_1=0.44/0.64=0.6875$，$b_2=-0.20/0.64=-0.3125$，$R^2=0.34375-0.03125=0.3125>0.25=r_{y1}^2$．

\subsection{混合分布の平均・分散}\label{subsec:G-mixture}
確率 $\pi_k$ で成分 $k$（平均 $\mu_k$，分散 $\sigma_k^2$）から生じる $X$ は，成分を表す確率変数 $C$ で条件付けると
\[
\E[X]=\sum_k\pi_k\mu_k,\qquad \V[X]=\E[\V[X\mid C]]+\V[\E[X\mid C]]=\sum_k\pi_k\sigma_k^2+\sum_k\pi_k(\mu_k-\E[X])^2
\]
（群内分散の平均＋群平均の分散）．2群の要約統計量から全体の平均・SD を求める計算もこの分解による．
混合率 $1/2$，等分散 $\sigma^2$ の正規混合の密度が二峰になるのは $|\mu_1-\mu_2|>2\sigma$ のときである（中点 $\xi=(\mu_1+\mu_2)/2$ で $f'(\xi)=0$ となり，$f''(\xi)>0$ なら $\xi$ は極小）．

\subsection{誤分類・ランダム化回答と境界での最尤推定}\label{subsec:G-rr}
真の割合 $\pi$ の属性を，誤回答率（偽陰性 $h_1$，偽陽性 $h_2$）のある方法で調べると，「はい」の確率は全確率の公式で $q=(1-h_1)\pi+h_2(1-\pi)$ である（診断検査の感度 $1-h_1$，偽陽性率 $h_2$ と同じ構造：\cref{ch:05}）．
はいの数 $m\sim\Bin(n,q)$ なので，$q$ の MLE は $m/n$ で，$q$ と $\pi$ が1対1（$1-h_1-h_2\ne0$）なら MLE の不変性から
$\hat\pi=(m/n-h_2)/(1-h_1-h_2)$．
\textbf{境界}：$\pi\in[0,1]$ に対応する $q$ の範囲は $[h_2,1-h_1]$（$1-h_1>h_2$ のとき）なので，$m/n$ がこの外にあれば尤度は範囲の端で最大になり，$\hat\pi=0$ または $1$ とする．
ランダム化回答法（確率 $p$ で「属性あり？」，$1-p$ で「属性なし？」と尋ねる）では $q=p\pi+(1-p)(1-\pi)$，$\hat\pi=\{m/n-(1-p)\}/(2p-1)$，
$\V[\hat\pi]=q(1-q)/\{n(2p-1)^2\}$ で，$p$ が $1/2$ に近いほどプライバシーは守られるが分散は大きい．

\subsection{離散分布の正確な棄却域と検出力}\label{subsec:G-exact}
離散分布では，有意水準 $\alpha$ を超えない最大の棄却域を選ぶので，実際の第1種の過誤確率は $\alpha$ より小さくなりうる．
例：$X\sim\Bin(10,p)$，$H_0:p=1/2$ を $H_1:p>1/2$ に対して $\alpha=0.05$ で検定する．
$\Prob_0(X\ge9)=11/1024=0.011$，$\Prob_0(X\ge8)=56/1024=0.055>0.05$ なので棄却域は $X\ge9$．
$p=0.8$ での検出力は $\Prob(X\ge9)=10\times0.8^9\times0.2+0.8^{10}=0.268+0.107=0.376$．
正規近似（連続修正の有無）で棄却域を決めると，正確な計算と境界が変わることがあるので，問題文の指示に従う．

\subsection{平均への回帰}\label{subsec:G-rtm}
$(X,Y)$ が平均 $\mu$，分散 $\sigma^2$，相関 $\rho$ の2変量正規なら $\E[Y\mid X=x]=\mu+\rho(x-\mu)$．
$x$ が平均より高い人を選ぶと，治療しなくても次の測定値は平均 $(1-\rho)(x-\mu)$ だけ下がる．
真値＋測定誤差のモデル $X=\theta+\varepsilon_1$，$Y=\theta+\varepsilon_2$（$\theta\sim\Normal(\mu,\tau^2)$，$\varepsilon_j\sim\Normal(0,\psi^2)$，独立）では $\rho=\tau^2/(\tau^2+\psi^2)$ で，
測定誤差が大きいほど回帰は強い．$\theta\mid X$ の分布は正規--正規モデルの事後分布（\cref{sec:17-normal}）と同じ計算である．
対照群をおけば平均への回帰は両群に同じく働くので，群間比較（\cref{sec:08-ancova}）で取り除ける．

\section{「着手しやすい問題」と「計算量が多い問題」の観察的特徴}\label{sec:G-features}
9回の問5を，観察できる特徴で比べた．ここで「着手しやすい」は難易度ではなく，最初の数分で解答の見通しが立つかどうかを指す．
\begin{center}\small
\begin{tabular}{@{}L{25mm}L{55mm}L{58mm}@{}}
\toprule
観点 & 着手しやすい側に見られた特徴 & 計算量が多い側に見られた特徴\\
\midrule
式変形の段数 & 1〜3段（代入，標準化，2×2 の逆行列） & 5段以上（正規密度の2回微分，共分散行列の固有分解，経路の全列挙）\\
数値計算 & 数表1〜2回＋四則演算 & 平方根を含む2次方程式，分数の積和，累乗\\
誘導 & 公式・係数が問題文にある & 誘導がない，または式変形の一部だけ\\
小問の独立性 & 後半が前半の結果を使わない & 後半が前半の式を直接使う\\
落とし穴 & 少ない & 丸めで答えが変わる，場合分け（境界解，符号），厳密性（極値の個数）\\
記述 & 1〜2文で書ける & 「詳細に述べよ」と長い説明を求める\\
\bottomrule
\end{tabular}
\end{center}
どの年も\textbf{最初の小問は定義の確認か公式への代入}（2018年はヒント付きの短い証明）であり，末尾は記述・一般化で終わる年が多い．
したがって問5は，3問目として選び，前半の小問だけで途中点を確保するという使い方も考えられる（解答する3問の選び方と答案の書き方は，当日の問題冊子の指示に従うこと）．

\section{他者の経験（事実と意見の区別）}\label{sec:G-cases}
本書の改訂にあたり，2026年10月24日〜11月1日に収集された合格体験記・SNS 投稿・GitHub 上の資料の調査記録（\cite{web:others}）を参照した．
医薬生物学に関する記事は少なく（調査記録では全体で10本前後），以下は個人の申告であって検証されたものではない．
\begin{center}\small
\begin{tabular}{@{}L{32mm}L{40mm}L{66mm}@{}}
\toprule
事例（年・分野） & 問5の扱い & 結果・本人の評価（区別）\\
\midrule
2019年・応用（分野不明） & 選択 & 【申告】基礎的な問5を完答し，他の大問は小問1を拾って優秀成績賞\\
2023年・医薬 & 「必ず選ぶ」と決めていた & 【申告】モンティ・ホール問題で想定が外れ苦戦（合否の記載なし）\\
2025年・医薬 & 選択せず（2・3・4） & 【申告】自己採点5割で不合格．【意見】「2問をじっくり解いて問5〔1〕だけ足せば十分だった」\\
2025年・医薬（2名） & 1名は選択（2・3・5），1名は選択せず（1・2・3） & 【申告】結果の記載なし\\
年・分野不明（合格者） & 避ける戦略 & 【意見】過去問をジャンル別に分類し，得意な形式で選ぶ\\
解説ブログ（2014〜2021年の分析） & --- & 【意見】「共通問題は医薬の応用問題より解きやすいものが多い」（根拠は1件のみ）\\
\bottomrule
\end{tabular}
\end{center}
\paragraph{医薬生物学全体について他者が述べていること}
\begin{itemize}
\item 【申告】2025年の医薬では，問2でポアソン分布のガンマ混合（負の二項分布）が出たこと（本書\cref{sec:17-gammapois}），順位和検定が出たことが記録されている．公式問題で確認した内容（\cref{app:D}）と一致する．
\item 【意見】2025年について「数理統計の力が直接問われた」「数理統計がきちんとできないと解けなくなってきている」という声が複数ある．2024年に医薬で優秀成績賞を得た受験者も「まず統計数理の対策を」と書いている．
\item 【意見】「迷う場面が3回以上あると時間切れになる」「解けない小問に早く見切りをつければ得点が伸びた」（他分野・数理の受験者）．
\item 【申告】対策した分野（分割表・ノンパラ・生存）が「どれも出なかった」年（2023年）の記録がある．なお2023年の医薬問2は打ち切りのない指数分布の最尤推定と症例数設計であり（\cref{app:D}），「生存時間が数理寄りの形で出た」という別の受験者の記録と矛盾しない．
\item 【宣伝・予想】有料教材や講座の目次（MH・SMR，RMST，Fine--Gray，MMRM・LOCF など）や出題予想は，出題実績ではないので本書の改訂の根拠にはしていない．
\end{itemize}

\section{試験本番の手順（本書の提案）}\label{sec:G-plan}
以下は過去問の構造と上の事例から本書が提案する手順であり，唯一の正解ではない．過去問演習で自分に合う形に調整する．

\subsection{最初の数分：3問の選び方}
\begin{enumerate}
\item \textbf{全5問の〔1〕〔2〕と最終小問だけを読む}（1問1分程度）．問題文の表・式の量にひるまず，「何を求めるか」だけを拾う．
\item 各問について次の4点を書き込む（問題用紙の余白に記号で）．
  \begin{itemize}
  \item[(a)] 既知のテーマか（本書のどの章か）．医薬固有の問1〜4は，準備した分野が出たかどうかで判断できる．
  \item[(b)] 誘導に乗れるか：与えられた式を見て「尤度を書く」「期待値・分散を計算する」「代入する」のどれかに落とせるか．
  \item[(c)] 計算量：\cref{sec:G-features}の右列の特徴（長い式変形，2次方程式，全列挙，場合分け）があるか．電卓での数値計算が長い表（多数の層・時点）があるか．
  \item[(d)] 小問の独立性：前半で詰まっても後半を拾えるか．
  \end{itemize}
\item (a)〜(d) で見通しの立つ2問を確定し，3問目は残りから「(b) と (d) が良いもの」を選ぶ．問5もこの基準で他の問と同列に比べる．
\item 選択に5分以上かけない．迷ったら，途中点を取りやすい（小問が独立で，前半が定義・代入の）問を選ぶ．
\end{enumerate}

\subsection{医薬固有問題（問1〜4）と問5の比較}
\begin{center}\small
\begin{tabular}{@{}L{32mm}L{53mm}L{53mm}@{}}
\toprule
観点 & 問1〜4（医薬固有） & 問5（共通）\\
\midrule
前提知識 & 医薬固有の手法（生存時間，分割表，診断，メタアナリシス，多重性など）．誘導が付く設問も多い & 確率・推測の基礎で最後まで解ける年が多い\\
予測しやすさ & 分野が比較的絞られる（生存時間はほぼ毎回，分割表は形を変えて繰り返す） & 9回で同じテーマなし\\
解釈の小問 & 交絡・非可縮性・競合リスクなど医薬の文脈の知識が要る & 数式をほぼ使わない記述が多い\\
判断材料 & テーマ名と表の形で数秒で分かる & 〔1〕〔2〕を読めば後半の計算量を見積もれる\\
\bottomrule
\end{tabular}
\end{center}
\textbf{判断基準の例}：準備した医薬の分野が2問以上出ていれば，それを優先し，3問目を問5を含む残りから選ぶ．
準備した分野が1問しかなければ，問5の〔1〕〔2〕を読み，(b)(c)(d) が良ければ問5を2問目の候補にする．
どちらの場合も「問5だから選ぶ／捨てる」と事前に決めつけない．

\subsection{途中点の確保}
\begin{itemize}
\item \textbf{定義・モデルを最初に書く}：確率モデル（どの変数がどの分布に従うか），推定対象，帰無仮説を式で書く．これだけで点になる設問が多い．
\item \textbf{導出は「方針→式→結果」の順}：最後まで計算できなくても，尤度・スコア・分散の式を立てた段階を答案に残す．
\item \textbf{前の小問の結果が出なくても先へ進む}：「〔2〕の結果を $\hat\theta$ とおくと」と記号で進めれば，後の小問の考え方は採点されうる．
\item \textbf{数値は途中の値も書く}：重み，期待値，分散などの中間値を表にして書くと，計算ミスがあっても方針の正しさが伝わる．
\item \textbf{解釈の小問は必ず書く}：医薬では「推定対象の違い」「仮定」「臨床的な意味」を1〜3文で書く設問がほぼ毎回あり，計算が苦手でも得点源になる．
\end{itemize}

\subsection{90分の時間配分}
\begin{center}\small
\begin{tabular}{@{}lL{110mm}@{}}
\toprule
時間 & 内容\\
\midrule
0〜5分 & 全問の〔1〕〔2〕・最終小問を読み，3問を決める\\
5〜35分 & 1問目（最も見通しの立つ問）．25分経過で残りの小問の見通しを確認\\
35〜65分 & 2問目\\
65〜85分 & 3問目（解ける小問から．時間が足りなければ定義・モデル・方針を書く）\\
85〜90分 & 見直し：記号の定義，自由度，片側・両側，単位，対数スケールで統合した結果の指数変換\\
\bottomrule
\end{tabular}
\end{center}
\begin{itemize}
\item 1問に30分を超えそうになったら，その時点で残りの小問は方針だけ書いて次へ進む（他者の申告では，1問に時間をかけすぎたこと・迷う場面が重なったことが失敗の原因として繰り返し挙げられている）．
\item 3問とも完答を目指すより，2問を確実に仕上げ3問目で途中点を拾う方が安定する，という振り返りもある（【意見】）．自分の演習結果で判断する．
\item 過去問演習は本番と同じ90分・3問で行い，選択にかかった時間も記録する．
\end{itemize}
